\section{Unified Connectivity Construction}

\subsection{Graph and construction metadata}

We represent layer $l$ by $m_l$ two-input gates and a fixed predecessor matrix $I_l\in\{0,\ldots,n_l-1\}^{2\times m_l}$. Here, $n_l$ is the preceding layer's width, and gate $j$ receives distinct predecessors $a=I_l[0,j]$ and $b=I_l[1,j]$. During training, the gate uses the usual mixture of the 16 relaxed Boolean functions \cite{petersen2022}:

\begin{equation}
h_{l,j}=\sum_{q=0}^{15}\mathrm{softmax}(\theta_{l,j})_q\,\widetilde f_q(h_{l-1,a},h_{l-1,b}).
\label{eq:gate}
\end{equation}

Training optimizes $\theta$ while keeping $I_l$ fixed; hardening then selects the highest-weight Boolean function, and class-group summation produces the output scores. Our method changes only the construction of $I_l$; the loss, gate parameterization, and inference operators remain those of the corresponding baseline.

To construct these indices, we associate each predecessor $i$ with a set $A_i$ of structural ancestry labels. In a dense image network, an input label identifies one raw channel-pixel value, so multiple thermometer thresholds of that value share a label. Propagation takes the union $A_j=A_a\cup A_b$, recording potential structural access without claiming functional dependence after gate learning. Figure \ref{fig:overview} shows how the metadata supports construction and is then discarded.

\begin{figure*}[t]
\centering
\begin{tikzpicture}[x=1cm,y=1cm,>=Stealth,
 box/.style={draw=methodblue,rounded corners=1pt,align=center,minimum height=0.85cm,font=\normalsize},
 smallbox/.style={box,text width=3.5cm},
 arr/.style={->,thick,draw=methodblue}]
\node[smallbox] (dense) at (1.8,1.0) {Dense image input\\$(c,y,x,\text{threshold})$};
\node[smallbox] (conv) at (1.8,-0.25) {Convolutional channels\\retain spatial sampling};
\node[box,text width=3.25cm] (semantic) at (6.0,1.0) {Semantic pairing\\first layer / classifier input};
\node[box,text width=3.25cm] (stages) at (6.0,-0.25) {Multiscale matchings\\deeper layers};
\node[box,text width=3.4cm] (fixed) at (10.4,0.4) {Fixed predecessor indices\\train gate functions};
\node[box,text width=2.4cm] (hard) at (14.1,0.4) {Harden gates\\Boolean inference};
\draw[arr] (dense) -- (semantic);
\draw[arr] (conv.east) -- ++(0.35,0) |- (semantic.west);
\draw[arr] (semantic) -- (stages);
\draw[arr] (semantic.east) -- ++(0.4,0) |- (fixed.west);
\draw[arr] (stages.east) -- ++(0.4,0) |- (fixed.west);
\draw[arr] (fixed) -- (hard);
\node[align=center,font=\normalsize,text=controlgray] at (7.6,-1.2)
 {Stage priority: degree spread $\rightarrow$ ancestry novelty $\rightarrow$ nominal scale};
\end{tikzpicture}
\caption{Overview of LogicConnect. Semantic pairing constructs input connections, and multiscale matching organizes deeper connections. The convolutional construction preserves spatial sampling and shared logic trees. Training optimizes gate functions on the resulting fixed graph.}
\label{fig:overview}
\end{figure*}


\subsection{Semantic input pairing}

At the image input, ancestry alone does not distinguish useful pairings, so we retain each bit's channel, spatial coordinates, and threshold index. Candidate stages pair positions along the horizontal and vertical axes at increasing power-of-two strides, capped at half the axis length. The finest scale also includes aligned channel pairs and threshold-changing pairs with a spatial displacement, excluding pairs from the same raw source.

The constructor cycles through these semantic stages to represent several axes rather than exhausting one axis first. Within the selected stage, it examines up to 64 unused candidates and chooses the pair minimizing endpoint fan-out sum, then maximum endpoint fan-out, then endpoint indices. A seeded affine permutation determines candidate order without consuming the training generator. Once a stage's distinct pairs are exhausted, repeated pairs become eligible using the same degree ordering over the stage. This procedure provides a reproducible semantic input graph without claiming globally optimal balance.

\subsection{Deeper multiscale matchings}

For deeper layers, we replace semantic-coordinate stages with regular matchings over predecessor indices, using $K=\lceil\log_2 n\rceil$ scales for $n$ inputs. At power-of-two widths, scale $s$ pairs $i$ with $i\mathbin{\oplus}2^s$. Otherwise, an affine traversal with a coprime step visits each predecessor once, and adjacent elements form pairs. Odd widths leave one predecessor unmatched, with the omitted index rotating across the schedule. For example, eight predecessors yield four disjoint pairs at each of three XOR scales. Filling ten gates uses two complete matchings and two pairs from a third, giving predecessor degrees of two or three.

These matchings control endpoint reuse independently of scale selection: $d_i$ denotes current fan-out, and $P_s$ contains pairs retained from candidate stage $s$. If only part of a stage fits the remaining width, we select its disjoint pairs by endpoint degree sum, maximum degree, and index. We then evaluate the prospective spread

\begin{equation}
D_s=\max_i(d_i+\delta_i(P_s))-\min_i(d_i+\delta_i(P_s)),
\label{eq:degree}
\end{equation}

where $\delta_i(P_s)$ counts occurrences of predecessor $i$ in the retained pairs. Among stages attaining the minimum spread, the constructor favors lower ancestry overlap. For the complete candidate matching $M_s$, its score is

\begin{equation}
N_s=\frac{1}{|Q_s|}\sum_{(a,b)\in Q_s}\left(1-\frac{|A_a\cap A_b|}{\max(1,\min(|A_a|,|A_b|))}\right).
\label{eq:novelty}
\end{equation}

Here, $Q_s$ contains all pairs in $M_s$, or 2,048 uniformly indexed pairs for a larger stage. The score approaches zero for identical ancestry and one for disjoint ancestry, using the complete candidate stage even when only a subset fits. Equal scores follow the cyclic nominal scale order, and each scale is used once before a new cycle begins. Algorithm \ref{alg:construction} summarizes this selection.

\begin{algorithm}[t]
\caption{Deeper-layer stage selection}
\label{alg:construction}
\begin{algorithmic}
\REQUIRE Ancestry $A$, width $m$, layer $l$, seed
\STATE Initialize degrees $d=0$, output pairs $I=\varnothing$.
\WHILE{$|I|<m$}
  \STATE Generate the $K$ seeded scale matchings $M_s$.
  \STATE Form degree-ordered partial stages $P_s$ if needed.
  \STATE Compute $D_s$ and $N_s$ using (\ref{eq:degree})--(\ref{eq:novelty}).
  \STATE Among unused scales, minimize $D_s$, then maximize $N_s$; break ties by nominal scale order.
  \STATE Append $P_s$; update $d$ and mark scale $s$ used.
  \STATE Reset used scales after a complete $K$-stage cycle.
\ENDWHILE
\STATE Propagate ancestry by union; return fixed indices $I$.
\end{algorithmic}
\end{algorithm}


For even widths, every complete matching increments each predecessor's degree once, while a final subset increments a predecessor at most once. An initially uniform degree vector therefore ends with spread at most one after the complete stages and the final subset. This property follows from disjoint endpoints, but native dense random connectivity already balances fan-out and therefore provides an equally relevant control. A degree-preserving ablation is necessary to examine whether the particular predecessor combinations provide value beyond their degree distribution.

\subsection{Convolutional adaptation}

The convolutional adaptation applies the pair constructor to channels while retaining each logic tree's spatial receptive-field sampler, depth, and weight sharing. The evaluated architecture uses two-channel groups, depth-three trees, and OR pooling, with trainable gate functions shared across spatial positions \cite{petersen2024}. First-layer pairing uses channel and threshold semantics, while later layers use the multiscale matching rule to select their channel groups. Thus, the construction changes channel access without replacing convolution with a dense spatial graph.

For this channel representation, ancestry starts with encoded-channel identities and gains a distinct identity label for each subsequent output channel. At flattening, each feature retains its upstream labels and receives a spatial-feature identity. The classifier first uses semantic pairing over the resulting channel-position grid, then the deeper matching rule. These labels prevent channel-support saturation from erasing all distinctions, but remain construction metadata rather than exact pixel-dependence sets.

\subsection{Cost and deployment}

Across both families, construction uses packed ancestry bitsets on the CPU before training. Candidate-stage scoring adds offline work and temporary storage, whose construction cost must be distinguished from GPU allocation during training. For a raw dense layer with $m$ gates, the model retains $16m$ trainable gate logits and $2m$ fixed predecessor indices, with no trainable routing state. Export discards the ancestry metadata and retains the selected Boolean functions and indices; physical implementation costs still require synthesis and routing measurements.
